% subsubsection: 1次分数型  \, $\displaystyle a_{n+1}=\frac{pa_n+q}{ra_n+s}$

\subsubsection{\texorpdfstring{1次分数型  \, $\displaystyle a_{n+1}=\frac{pa_n+q}{ra_n+s}$}{1次分数型}}\label{節：1次分数型}
    先程の逆数型の一般化と言える漸化式である,\,1次分数型の漸化式について考えていこう.
    以下では,\, $r\ne0$,\, $ps-qr\ne0$とし,もとの漸化式の分母が$0$にならないものとして議論を進める.

    数学における基本的な指針として,既知の形に帰着させるというものがある.
    本書で言う,"きれい"な形にするということである.
    今回の漸化式で,分子の定数項$q$がなければ,つまりは,
    \[
      a_{n+1}=\frac{pa_n+q}{ra_n+s} \Leftrightarrow b_{n+1}=\frac{tb_n}{u b_n+v}
    \]
    となるような$b_n$と$t,u,v$の組があれば逆数型として解くことができる.
    そこで,分子の定数項を消し,逆数型に帰着させることを目標にしよう.

    \,\sectionref{節：特性方程式型}では,数列の各項から同じ定数を引くことで定数項を消した.
    今回もそのアイデアを使い,\, $b_n=a_n-\alpha$と置いてみよう.
    ただし,ここでは$\alpha$をまだ決めず,変形後の分子の定数項が消えるように選ぶ.
    $a_n=b_n+\alpha$を代入すると,
    \begin{align*}
      b_{n+1}-\alpha &=\frac{p(b_n+\alpha)+q}{r(b_n+\alpha)+s}\\
      b_{n+1} &=\frac{p(b_n+\alpha)+q}{r(b_n+\alpha)+s}-\alpha\\
      &=\frac{p(b_n+\alpha)+q-\alpha\{r(b_n+\alpha)+s\}}
        {r(b_n+\alpha)+s}\\
      &=\frac{(p-r\alpha)b_n+p\alpha+q-r\alpha^2-s\alpha}
        {rb_n+r\alpha+s}
    \end{align*}
    となる.したがって,
    \[
      p\alpha+q-r\alpha^2-s\alpha=0
    \]
    となるように$\alpha$を選べば,
    \[
      b_{n+1}=\frac{(p-r\alpha)b_n}{rb_n+r\alpha+s}
    \]
    という逆数型に帰着できる.
    定数項を消すための条件は,
    \[
      r\alpha^2+(s-p)\alpha-q=0
      \quad\Longleftrightarrow\quad
      \alpha=\frac{p\alpha+q}{r\alpha+s}
    \]
    と書けるので,この$\alpha$はもとの漸化式の不動点でもある.
    なお,\, $r\alpha+s=0$なら上の方程式から$p\alpha+q=0$も成り立ち,\, $ps-qr=0$となって仮定に反するので,この分母は$0$ではない.

    ここまでで逆数型として解くことができるが,定数項を消せる値を求める問題は,
    \[
      rx^2+(s-p)x-q=0
    \]
    という二次方程式を解く問題に帰着した.
    この方程式が異なる2つの実数解$\alpha,\beta$を持つなら,その左辺は
    \[
      rx^2+(s-p)x-q=r(x-\alpha)(x-\beta)
    \]
    と因数分解できる.
    $\alpha,\beta$は別々に選ぶ定数ではなく,もとの漸化式の係数から決まる,同じ二次方程式の二つの解なのである.
    そして,どちらの解も定数項を消す条件を満たすので,\, $a_n-\alpha$と$a_n-\beta$のどちらについても,同じ変形ができる.
    先ほどの逆数型の式を,もとの$a_n$を使って書き直すと,
    \begin{align*}
      a_{n+1}-\alpha
      &=\frac{p-r\alpha}{ra_n+s}(a_n-\alpha),\\
      a_{n+1}-\beta
      &=\frac{p-r\beta}{ra_n+s}(a_n-\beta)
    \end{align*}
    となる.

    この二つの式から,さらに単純な法則を取り出せないだろうか.
    それぞれの差に掛かる倍率には$a_n$が含まれているので,差そのものは一般には等比数列にならない.
    しかし,二つの倍率には同じ分母$ra_n+s$が現れている.
    そこで二つの式の比を取れば,この分母が消え,一定の倍率だけが残りそうである.

    初項が$\alpha$または$\beta$なら定数数列になるので,先に扱っておこう.
    以下では初項がどちらの不動点とも異なる場合を考える.
    また,\, $p-r\beta=0$なら不動点の条件から$q=s\beta$となり,\, $ps-qr=0$となって仮定に反する.
    よって,\, $p-r\beta\ne0$であり,上の式から$a_n\ne\beta$なら$a_{n+1}\ne\beta$である.
    初項から順に考えると,すべての項で$a_n\ne\beta$なので,二つの差の比を取ることができる.
    実際に,
    \begin{align*}
      \frac{a_{n+1}-\alpha}{a_{n+1}-\beta}
      &=\cfrac{\displaystyle
        \frac{(p-r\alpha)(a_n-\alpha)}{ra_n+s}}
        {\displaystyle
        \frac{(p-r\beta)(a_n-\beta)}{ra_n+s}}\\
      &=\frac{p-r\alpha}{p-r\beta}
        \frac{a_n-\alpha}{a_n-\beta}
    \end{align*}
    となる.
    したがって,二つの差の比は,項が一つ進むごとに一定の倍率$\dfrac{p-r\alpha}{p-r\beta}$で変化する.
    分子の定数項を消せる二つの式を比べることで,等比数列の法則を見つけられたのである.

    ここで
    \[
      b_n=\frac{a_n-\alpha}{a_n-\beta}
    \]
    と置けば,
    \[
      b_{n+1}
      =
      \frac{p-r\alpha}{p-r\beta}b_n
    \]
    となる.したがって,\,$b_n$は等比数列である.
    $\alpha$と$\beta$の大小関係は,解の大小関係を要求されていないことから自由に選んでよい.
    つまり,\,1次分数型漸化式では,
    \[
      a_{n+1}=\frac{pa_n+q}{ra_n+s}
    \]
    という複雑な形をそのまま扱うのではなく,
    \[
      \frac{a_n-\alpha}{a_n-\beta}
    \]
    という「$\alpha,\beta$からの距離の比」に置き換えることで,
    等比数列に変えることができる.

    そして,\,$\alpha,\beta$は,
    \[
      rx^2+(s-p)x-q=0
    \]
    の異なる2解として自然に現れる.
    
    したがって\,$\alpha,\beta$は,もとの漸化式の$a_n$を形式的に$\alpha,\beta$に置き換えたときにも値が変化しない数,すなわち不動点である.
    
    議論が長くなったので,ここで解法をまとめておこう.
    \[
      a_{n+1}=\frac{pa_n+q}{ra_n+s} 
    \]
    の形で表される漸化式に対して,ただし$ps-qr\ne0$とし,
    $r\ne0$とする.また,以下では$a_n\ne-\dfrac sr$(漸化式の分母が$0$にならないこと)を仮定する.
    さらに,$a_1\ne\beta$とする.
    \[
      x= \frac{px+q}{rx+s} \Leftrightarrow rx^2+(s-p)x-q=0
    \]
    を解き,その異なる2解を$\alpha,\beta$とする.
    ここで,
    \[
      b_n=\frac{a_n-\alpha}{a_n-\beta}
    \]
    と置く.
    すると,
    \begin{align*}
      b_{n+1}&=\frac{a_{n+1}-\alpha}{a_{n+1}-\beta}\\
      &=\cfrac{\displaystyle\frac{pa_n+q}{ra_n+s}-\alpha}{\displaystyle\frac{pa_n+q}{ra_n+s}-\beta}\\
      &=\frac{pa_n+q-\alpha(ra_n+s)}{pa_n+q-\beta(ra_n+s)}\\
      &=\frac{(p-r\alpha)a_n+q-s\alpha}{(p-r\beta)a_n+q-s\beta}
    \end{align*}
    ここで,\,$\alpha$は
    \[
      r\alpha^2+(s-p)\alpha-q=0
    \]
    を満たすから,
    \[
      q-s\alpha=-\alpha(p-r\alpha)
    \]
    である.
    したがって,
    \[
      (p-r\alpha)a_n+q-s\alpha
      =(p-r\alpha)(a_n-\alpha)
    \]
    同様に,\,$\beta$について
    \[
      (p-r\beta)a_n+q-s\beta
      =(p-r\beta)(a_n-\beta)
    \]
    である.
    よって,
    \begin{align*}
      b_{n+1}&=\frac{(p-r\alpha)(a_n-\alpha)}{(p-r\beta)(a_n-\beta)}\\
      &=\frac{p-r\alpha}{p-r\beta}\frac{a_n-\alpha}{a_n-\beta}\\
      &=\frac{p-r\alpha}{p-r\beta}b_n
    \end{align*}
    したがって,
    \[
      b_{n+1}=\frac{p-r\alpha}{p-r\beta}b_n
    \]
    となり,\,$b_n$は等比数列である.
    初項
    \[
      b_1=\frac{a_1-\alpha}{a_1-\beta}
    \]
    を用いれば,
    \[
      b_n=\frac{a_1-\alpha}{a_1-\beta}\left(\frac{p-r\alpha}{p-r\beta}\right)^{n-1} \tag*{\(\cdots\) (＊)}
    \]
    となる.
    最後に,
    \[
      b_n=\frac{a_n-\alpha}{a_n-\beta}
    \]
    から$a_n$を求める.
    \begin{align*}
      b_n(a_n-\beta)&=a_n-\alpha\\
      (b_n-1)a_n&=\beta b_n-\alpha\\
      a_n=\frac{\beta b_n-\alpha}{b_n-1}
    \end{align*}
    (＊)を代入して$a_n$を求めることができる.

    なお,$rx^2+(s-p)x-q=0$が重解$\alpha$を持つ場合は,$\beta$と異なる2つの不動点を取れないため,上の議論はそのまま使えない.
    このときは$\alpha=\dfrac{p-s}{2r}$,すなわち$s=p-2r\alpha$が成り立つため,
    \[
      c=p-r\alpha=r\alpha+s\ne0
    \]
    とおくと,
    \[
      a_{n+1}-\alpha
      =\frac{c(a_n-\alpha)}{r(a_n-\alpha)+c}
    \]
    となる.したがって,
    \[
      \frac{1}{a_{n+1}-\alpha}
      =\frac{1}{a_n-\alpha}+\frac{r}{c}
    \]
    であるから,右辺も含めて$\displaystyle \frac{1}{a_n-\alpha}$について考えれば等差数列に帰着できる.
    ここで$b_n=a_n-\alpha$とおくと,
    \[
      b_{n+1}=\frac{cb_n}{rb_n+c}
    \]
    という逆数型漸化式に帰着する(先ほどの一般式で$\beta=\alpha$とした場合に対応する).したがって両辺の逆数を取れば
    \[
      \frac1{b_{n+1}}=\frac1{b_n}+\frac rc
    \]
    という等差数列が得られ,ここから$a_n$を求めることができる(具体的な計算は例題(3)を参照).

    \noindent \bookblocklink{example-mobius-1}{problem}{answer}{【例題】}
    \begin{enumerate}[label=(\arabic*),parsep=3pt,format=\bookitemlink{example-mobius-1}{problem}{answer}]
      \item $\displaystyle a_1=3,\, a_{n+1}=\frac{2a_n+1}{a_n+2}$
      \item $\displaystyle a_1=0,\, a_{n+1}=\frac{a_n+4}{a_n+1}$
      \item $\displaystyle a_1=5,\, a_{n+1}=\dfrac{4a_{n}-9}{a_{n}-2}$
      \item $\displaystyle a_1=1,\, a_{n+1}=1+\frac{1}{a_n}$
    \end{enumerate}

    \noindent\bookblocklink{example-mobius-1}{hint}{problem}{【方針】}
    \begin{enumerate}[label=(\arabic*),parsep=3pt,format=\bookitemlink{example-mobius-1}{hint}{problem}]
      \item 最も基本的な例です.先程の議論に倣って解きましょう.
      \item 同様に不動点$\alpha,\beta$を求めましょう.今回は片方が負の値になります.
      \item 今回は不動点の方程式が重解になります.\,$\dfrac{a_n-\alpha}{a_n-\beta}$は使えないので,\,$b_n=a_n-\alpha$と置いてから逆数型に帰着させる別の手法を使います.
      \item このような形も一次分数型です.解と係数の関係を活用して解きましょう.
    \end{enumerate}

    \noindent\bookblocklink{example-mobius-1}{answer}{problem}{【解答】}
    \begin{enumerate}[label=(\arabic*),parsep=3pt,format=\bookitemlink{example-mobius-1}{answer}{problem}]
      \item $\displaystyle a_n=\frac{2\cdot3^{n-1}+1}{2\cdot3^{n-1}-1}$
      \item $\displaystyle a_n=\frac{2\left(3^{n-1}+(-1)^n\right)}{3^{n-1}-(-1)^n}$
      \item $a_{n}=\dfrac{6n-1}{2n-1}$
      \item $\displaystyle a_n=\cfrac{\left(\cfrac{1-\sqrt{5}}{2}\right)^{n+1}-\left(\cfrac{1+\sqrt{5}}{2}\right)^{n+1}}{\left(\cfrac{1-\sqrt{5}}{2}\right)^n-\left(\cfrac{1+\sqrt{5}}{2}\right)^n}$
    \end{enumerate}

    \noindent\bookblocklink{example-mobius-1}{solution}{problem}{【解法】}
    \begin{enumerate}[label=(\arabic*),parsep=3pt,format=\bookitemlink{example-mobius-1}{solution}{problem}]
      \item $\displaystyle a_1=3,\, a_{n+1}=\frac{2a_n+1}{a_n+2}$\\
            不動点の方程式は
            \[
              x=\frac{2x+1}{x+2}
            \]
            であるから,その解をそれぞれ$\alpha,\beta$とすると\,$\alpha=1,\ \beta=-1$と取れる.
            このとき,
            \[
              b_n=\frac{a_n-1}{a_n+1}
            \]
            と置くと,
            \begin{align*}
              b_{n+1}&=\frac{a_{n+1}-1}{a_{n+1}+1}\\
              &=\cfrac{\cfrac{2a_n+1}{a_n+2}-1}{\cfrac{2a_n+1}{a_n+2}+1}\\
              &=\frac{2a_n+1-(a_n+2)}{2a_n+1+(a_n+2)}\\
              &=\frac{a_n-1}{3a_n+3}\\
              &=\frac{1}{3}\cdot \frac{a_n-1}{a_n+1}
            \end{align*}
            $\displaystyle b_1=\frac{a_1-1}{a_1+1}=\frac{2}{4}=\frac{1}{2}$から,
            \begin{align*}
              b_{n}=\frac{1}{2}\cdot\left(\frac{1}{3}\right)^{n-1}\\
              \therefore \quad b_n=\frac{1}{2\cdot3^{n-1}}
            \end{align*}
            最後に
            \[
              b_n=\frac{a_n-1}{a_n+1}
            \]
            を$a_n$について解くと,
            \begin{align*}
              &b_n(a_n+1)=a_n-1\\
              &(b_n-1)a_n=-1-b_n\\
              &a_n=\frac{1+b_n}{1-b_n}
            \end{align*}
            \[
              b_n=\frac{1}{2\cdot3^{n-1}}
            \]
            を代入して整理すれば,
            \[
              a_n=\frac{2\cdot3^{n-1}+1}{2\cdot3^{n-1}-1}
            \]
            を得る.
      \item $\displaystyle a_1=0,\, a_{n+1}=\frac{a_n+4}{a_n+1}$\\
            不動点の方程式は
            \[
              x=\frac{x+4}{x+1}
            \]
            であるから,その解をそれぞれ$\alpha,\beta$とすると\,$\alpha=2,\ \beta=-2$と取れる.
            このとき,
            \[
              \frac{p-r\alpha}{p-r\beta}=\frac{1-2}{1+2}=-\frac13
            \]
            なので,
            \[
              b_n=\frac{a_n-2}{a_n+2}
            \]
            と置くと,
            \begin{align*}
              &b_{n+1}=-\frac13 b_n,\quad b_1=\frac{a_1-2}{a_1+2}=\frac{-2}{2}=-1\\
              &\therefore \quad b_n=(-1)\left(-\frac13\right)^{n-1}=\frac{(-1)^n}{3^{n-1}}
            \end{align*}
            \[
              b_n=\frac{a_n-2}{a_n+2}
            \]
            を$a_n$について解くと,
            \[
              a_n=\frac{2(1+b_n)}{1-b_n}
            \]
            となるので,
            \[
              b_n=\frac{(-1)^n}{3^{n-1}}
            \]
            を代入し,分母分子に$3^{n-1}$を掛けて整理すると,
            \[
              a_n=\frac{2\left(3^{n-1}+(-1)^n\right)}{3^{n-1}-(-1)^n}
            \]
            を得る.
      \item $\displaystyle a_1=5,\, a_{n+1}=\dfrac{4a_{n}-9}{a_{n}-2}$
            不動点の方程式は
            \[
              x=\frac{4x-9}{x-2}
            \]
            であるから,その解は3.
            \[
            b_n=a_n-3
            \]
            とおくと,
            \begin{align*}
              b_{n+1}&=a_{n+1}-3\\
              &=\dfrac{4a_{n}-9}{a_{n}-2}-3\\
              &=\dfrac{a_{n}-3}{a_{n}-2}\\
              &=\dfrac{b_{n}}{b_{n}+1}
            \end{align*}
            逆数を取って
            \begin{align*}
              \frac{1}{b_{n+1}}&=\dfrac{b_{n}+1}{b_{n}}
            \end{align*}
            $b_1=a_1-3=2$,すなわち$\dfrac1{b_1}=\dfrac12$より,
            \begin{align*}
              &\frac{1}{b_n}=\frac{1}{2}+(n-1)\cdot 1\\
              &\frac{1}{b_n}=\frac{2n-1}{2}\\
              &b_n=\frac{2}{2n-1}\\
              &a_n-3=\frac{2}{2n-1}\\
              &a_n=\frac{6n-1}{2n-1}
            \end{align*}
      \item $\displaystyle a_1=1,\, a_{n+1}=1+\frac{1}{a_n}$\\
            式を変形して
            $\displaystyle a_1=1,\, a_{n+1}=\frac{a_n+1}{a_n}$
            不動点の方程式は
            \[
              x=\frac{x+1}{x}\tag*{\(\cdots\) (1)}
            \]
            であるから,その解をそれぞれ$\alpha,\beta$とすると
            \[
              \alpha=\frac{1+\sqrt{5}}{2},\ \beta=\frac{1-\sqrt{5}}{2}
            \]
            ここで(1)について,その解$\alpha,\beta$では
            \[
              \alpha+\beta=1,\, \quad \alpha\beta=-1\tag*{\(\cdots\) (2)}
            \]
            が解と係数の関係から成り立つ.したがって,
            \begin{align*}
              b_{n+1}&=\frac{a_{n+1}-\alpha}{a_{n+1}-\beta}\\
              &=\cfrac{\cfrac{a_n+1}{a_n}-\alpha}{\cfrac{a_n+1}{a_n}-\beta}\\
              &=\frac{a_n+1-\alpha a_n}{a_n+1-\beta a_n}\\
              &=\frac{(1-\alpha)a_n+1}{(1-\beta)a_n+1}\tag*{\(\cdots\) (3)}
            \end{align*}
            (2)より,(3)は
            \begin{align*}
              \frac{(1-\alpha)a_n+1}{(1-\beta)a_n+1}&=\frac{\beta a_n+1}{\alpha a_n+1}\\
              &=\frac{\beta a_n-\alpha\beta}{\alpha a_n-\alpha\beta}\\
              &=\frac{\beta (a_n-\alpha)}{\alpha (a_n-\beta)}\\
              &=\frac{\beta}{\alpha}b_n
            \end{align*}
            よって,
            \[
              b_{n+1}=\frac{\beta}{\alpha}b_n
            \]
            初項$b_1$は,
            \begin{align*}
              b_1=\frac{a_1-\alpha}{a_1-\beta}=\frac{1-\alpha}{1-\beta}=\frac{\beta}{\alpha}
            \end{align*}
            したがって,
            \[
              b_n=\left(\frac{\beta}{\alpha}\right)^n
            \]
            次に
            \[
              b_{n}=\frac{a_{n}-\alpha}{a_{n}-\beta}
            \]
            を$a_n$について解くと,
            \[
              a_{n}=\frac{\beta b_{n}-\alpha}{b_{n}-1}
            \]
            したがって,
            \begin{align*}
              a_{n}&=\cfrac{\beta \left(\cfrac{\beta}{\alpha}\right)^n-\alpha}{\left(\cfrac{\beta}{\alpha}\right)^n-1}  \\
              &=\frac{\beta^{n+1}-\alpha^{n+1}}{\beta^n-\alpha^n}\\
              &=\cfrac{\left(\cfrac{1-\sqrt{5}}{2}\right)^{n+1}-\left(\cfrac{1+\sqrt{5}}{2}\right)^{n+1}}{\left(\cfrac{1-\sqrt{5}}{2}\right)^n-\left(\cfrac{1+\sqrt{5}}{2}\right)^n}
            \end{align*}        
    \end{enumerate}
  余談にはなるが,最後の問題はフィボナッチ数列の2項間の比の極限値は黄金比$\dfrac{1+\sqrt{5}}{2}$に一致する.
  実際,\,$a_n=\dfrac{F_{n+1}}{F_n}$(ただし$F_n$は$F_1=F_2=1$から定まるフィボナッチ数列)であることが,同じ漸化式$a_{n+1}=1+\dfrac1{a_n}$から確かめられる.
  また,ここから派生したリュカ数列も同じ特性方程式$x^2=x+1$を持ち,隣接項の比が同じ黄金比に収束する.
  一方,同様に漸化式から派生する数列にトリボナッチ数列があるが,こちらは特性方程式が3次$(x^3=x^2+x+1)$になり,隣接項の比の極限値(トリボナッチ定数,約$1.839$)は黄金比とは異なる.
  これらは本書の趣旨から外れてしまうため割愛する.

  加えて,このタイプの漸化式は帰納法によって示すことが比較的容易なパターンが存在する.
  例えば(3)は解答の一般項から分かるように,分母分子それぞれが一次関数の形で書かれているため,想像しやすい.
  $a_1$から$a_3$ぐらいまでを求めさせてから帰納法で示させるような誘導問題が出た時は帰納法を使って解くことも視野に入れておくといいだろう.
  以下にその例を示す.\\
    \noindent\bookblocklink{example-mobius-2}{problem}{answer}{【例題】}
    \begin{quote}
      以下の漸化式
      \[
        a_1=5,\, a_{n+1}=\dfrac{4a_{n}-9}{a_{n}-2}
      \]
      を満たす数列$a_n$の一般項を帰納法を用いて求めよ.
    \end{quote}
    \bookblocklink{example-mobius-2}{hint}{problem}{【方針】}
    \begin{quote}
      実験的に代入計算をしてみると,
      \[
        a_1=\frac{5}{1},\,a_2=\frac{11}{3},\,a_3=\frac{17}{5},\,a_4=\frac{23}{7},\,a_5=\frac{29}{9}
      \]
      となる.分母が2ずつ,分子が6づつ増えていることから,一般項は
      \[
        a_n=\frac{6n-1}{2n-1}
      \]
      と予想する.
    \end{quote}
    \bookblocklink{example-mobius-2}{answer}{problem}{【解答】}
    \begin{quote}
      一般項$a_n$が
      \[
        a_n=\frac{6n-1}{2n-1}
      \]
      であることを示す.\\
      $n=1$の時,
      \[
        \frac{6-1}{2-1}=5
      \]
      よって成立する.\\
      ある正の整数$k$について,$n=k$が成り立つと仮定する.このとき,
      \[
        n=k+1
      \]
      のとき,
      \begin{align*}
        a_{k+1}&=\frac{4a_{k}-9}{a_{k}-2}\\
        &=\cfrac{4\cdot \cfrac{6k-1}{2k-1}-9}{\cfrac{6k-1}{2k-1}-2}\\
        &=\frac{4(6k-1)-9(2k-1)}{6k-1-2(2k-1)}\\
        &=\frac{6k+5}{2k+1}\\
        &=\frac{6(k+1)-1}{2(k+1)-1}
      \end{align*}
      したがって,$n=k+1$でも成立する.
      以上より
      \[
        a_n=\frac{6n-1}{2n-1}
      \]
      が示された.\qed

  \end{quote}

